\documentclass{article}
\usepackage[margin=28mm]{geometry}
\usepackage{amsmath,amssymb}
\usepackage{xurl}
\usepackage[hidelinks]{hyperref}
% Shared commands for notes in docs/paper-gaps/.

\DeclareUnicodeCharacter{2080}{\ifmmode{}_{0}\else\textsubscript{0}\fi}
\DeclareUnicodeCharacter{2081}{\ifmmode{}_{1}\else\textsubscript{1}\fi}
\DeclareUnicodeCharacter{2082}{\ifmmode{}_{2}\else\textsubscript{2}\fi}
\DeclareUnicodeCharacter{2099}{\ifmmode{}_{n}\else\textsubscript{n}\fi}

\newcommand{\Lean}{\textsc{Lean}}
\ExplSyntaxOn
\NewDocumentCommand{\leanid}{m}{
  \ifmmode
    \textsf{\detokenize{#1}}
  \else
    \group_begin:
    \urlstyle{sf}
    \tl_set:Nn \l_tmpa_tl {#1}
    \tl_replace_all:Nnn \l_tmpa_tl {\_} {_}
    \exp_args:NV \path \l_tmpa_tl
    \group_end:
  \fi
}
\ExplSyntaxOff
\providecommand{\tr}{\operatorname{tr}}
\newcommand{\tnleanIssueBase}{https://github.com/LionSR/TNLean/issues/}
\newcommand{\tnleanPrBase}{https://github.com/LionSR/TNLean/pull/}
\newcommand{\tnleanDocBase}{https://sirui-lu.com/TNLean/docs/}
\newcommand{\ghissue}[1]{\href{\tnleanIssueBase#1}{GitHub issue \##1}}
\newcommand{\ghpr}[1]{\href{\tnleanPrBase#1}{PR \##1}}
\newcommand{\doclink}[1]{\href{\tnleanDocBase#1.html}{\path{#1}}}

% Dirac notation and the identity operator, as in blueprint/src/macros/common.tex.
% \providecommand keeps note-local definitions authoritative.
\usepackage{dsfont}
\providecommand{\ket}[1]{|#1\rangle}
\providecommand{\bra}[1]{\langle #1|}
\providecommand{\braket}[2]{\langle #1 | #2 \rangle}
\providecommand{\ketbra}[2]{|#1\rangle\!\langle #2|}
\providecommand{\Id}{\mathds{1}}
\providecommand{\one}{\mathds{1}}
\providecommand{\C}{\mathbb{C}}
\providecommand{\MN}[1]{M_{#1}(\C)}
\providecommand{\MD}{M_D(\C)}

% Verdict marker: \gapnote{<kind>}{<status>}. Kind is the policy
% classification (clarification, local-correction, scope-restriction,
% unfaithful, false-source, open-gap); status is open, wip (elimination
% underway), resolved, or historical. Machine-read by the site index;
% prints a verdict line.
\newcommand{\gapnote}[2]{\par\noindent\textbf{Verdict:} #1 (#2).\par}

\providecommand{\gapnote}[2]{\par\noindent\textbf{Verdict:} #1 (#2).\par}
\title{The Truncated Dynamics in Hastings' Approximation $M_i$}
\author{}
\date{}
\begin{document}
\begin{center}
 {\Large The Truncated Dynamics in Hastings' Approximation $M_i$}
\end{center}
\medskip
\gapnote{local-correction}{resolved}

\section{The printed estimate}
Hastings, \emph{Solving gapped Hamiltonians locally}~\cite{Hastings2006Solving},
\href{https://arxiv.org/abs/cond-mat/0508554v3}{arXiv:cond-mat/0508554v3},
source \path{proj3new.tex}, lines 285--320, approximates the filtered term
$\tilde H_i^0$ by
\begin{align}
 M_i=\frac{\Delta E}{\sqrt{2\pi q}}\int dt\,e^{-(t\Delta E)^2/(2q)}\,
 e^{iH_{\rm loc}t}H_ie^{-iH_{\rm loc}t},\qquad
 H_{\rm loc}=\sum_{j:\,d(i,j)\le l_{proj}-R}H_j,
 \label{eq:hastings05_truncated_dynamics_mi}
\end{align}
and states, eq.~(apb),
\begin{align}
 \|\tilde H_i^0-M_i\|\le J\bigl(N(l_{proj})\,g(c_1l_{proj},l_{proj})
 +2e^{-(c_1l_{proj}\Delta E)^2/(2q)}\bigr),
 \label{eq:hastings05_truncated_dynamics_apb}
\end{align}
with $N(l_{proj})$ the number of sites $j$ with
$l_{proj}-R<d(i,j)\le l_{proj}+R$. Eq.~(mid) prints the weight as
$e^{-(t\Delta E)^2/2}$; the weight of eq.~(thd), $e^{-(t\Delta E)^2/(2q)}$, is
meant, since otherwise $M_i$ would not be normalized and the bound
(\ref{eq:hastings05_truncated_dynamics_apb}) would not follow from the proof at
lines 1458--1533.

\section{The proof and the formalized estimate}
The proof (lines 1469--1505) compares the two evolutions of $H_i$. With
$V=H-H_{\rm loc}$ and $\tau_t$, $\tau^{\rm loc}_t$ the two Heisenberg evolutions,
\begin{align}
 \|\tau_t(W)-\tau^{\rm loc}_t(W)\|\le\Bigl|\int_0^t
 \|[V,\tau^{\rm loc}_u(W)]\|\,du\Bigr|.
 \label{eq:hastings05_truncated_dynamics_duhamel}
\end{align}
The terms of $V$ whose support does not meet the support $K$ of
$\tau^{\rm loc}_u(W)$ commute with it, and the others are bounded by the
Lieb--Robinson bound for the truncated dynamics. The source then asserts that
$JS\int_0^tg(t',l_{proj}-2R)\,dt'$ is bounded by $g(t,l_{proj})$, citing
[lsmh] (line 1496).

The formalization differs in three points, none of which changes the order of
the estimate.
\begin{enumerate}
 \item The truncation is stated for any family $s$ of terms acting inside a set
  $K$ that contains the support $S$ of $W$. The shell consists of the omitted
  terms whose supports meet $K$, and $m$ is a lower bound on the graph
  distance from $S$ to each of them. With $K$ the sites within distance
  $l_{proj}$ of $i$ and $s$ the terms with $d(i,j)\le l_{proj}-R$, each shell
  term has $l_{proj}-R<d(i,j)\le l_{proj}+R$, so the shell count is at most
  $N(l_{proj})$, and $m$ may be taken as $l_{proj}-3R$ when $H_i$ acts within
  distance $R$ of $i$.
 \item The time integral in (\ref{eq:hastings05_truncated_dynamics_duhamel})
  is bounded on $|t|\le T$ by $T$ times the largest integrand, instead of the
  absorption into $g(t,l_{proj})$. The Lieb--Robinson bound for a shell term
  acting on $X_j$ is obtained from the single-site bound by the partial trace
  over $X_j$, which costs a factor $2|X_j|\le2v_R$. The resulting bound is
  \begin{align}
   \|\tilde H_i^0-M_i\|\le2v_RC|S|\,\|W\|\,J\,N\,T\,e^{vT-\mu m}
   +2\|W\|e^{-(T\Delta E)^2/(2q)},
  \end{align}
  for every $T\ge0$ and $\mu\ge0$, where $C=2+4|b_0|e^{2\mu R}$ and
  $v=2v_R|b_0|e^{2\mu R}+1$ are the constants of the graph Lieb--Robinson
  bound, $N$ is the shell count and $J$ bounds the norms of the terms. With
  $\mu>0$ fixed, $m=l_{proj}-3R$ and $T=c_1l_{proj}$ for a constant
  $0<c_1<\mu/v$, the exponent is $-(\mu-vc_1)l_{proj}+3\mu R$, so the first
  term is exponentially small in $l_{proj}$, as in
  (\ref{eq:hastings05_truncated_dynamics_apb}).
 \item The constants are explicit and are not identified with the source's
  $g$ and $c_1$; the parameters $q$ and $T$ are left free.
\end{enumerate}
The support of $M_i$, the bound $\|M_i\|\le\|H_i\|$ and the hermiticity of
$M_i$ hold as stated in the source.

\section{The exponential form of eq.~(facld)}
Let $H=\sum_k h_k$ have ground energy $E_0$ and a positive spectral gap
$\Delta E$. Write $P_{high}$ for the orthogonal projection onto the complement
of its ground eigenspace and $\Psi_a$ for a normalized ground eigenvector.
The source states eq.~(facld), $|P_{high}M_i\Psi_a\rangle|\le
J\,O(e^{-l_{proj}/l_{fac}})$, only to order of magnitude. Its metric and local
range hypotheses at lines 169--190 do not imply polynomial growth of spheres
or of the omitted shell. The earlier shell-count argument obtains this form
under an additional hypothesis: at projection radius $l$, let $M_l$ be the
Gaussian-filtered approximation of $W$ with $q=c_1l\Delta E$, evolved under
the terms of a family $s_l$ supported inside $K_l\supseteq S$. Assume that
the number $N_l$ of omitted terms meeting $K_l$ is at most $K_g(l+1)^k$, for
fixed $K_g\ge0$ and $k\in\mathbb N$. For every ground vector $\psi$ of $H$,
with $\mu>0$, $0<c_1v<\mu$ and $m=l-3R$, it gives
\begin{align}
  |P_{high}M_l\psi\rangle|\le\|W\|\,(3+C'J)\,e^{-l/l_{fac}}\,\|\psi\|,
  \qquad
  l_{fac}=\frac{2}{\min(c_1\Delta E,\ \mu-c_1v)},
  \label{eq:hastings05_truncated_dynamics_facld_decay}
\end{align}
with $C'$ independent of $l$, using $(x+1)^kx\le(k+1)!\,e^{a}a^{-(k+1)}e^{ax}$
for $a=1/l_{fac}$ to absorb the polynomial factor. A subexponential bound
$N_l\le C_\varepsilon e^{\varepsilon l}$ would suffice with $l_{fac}$ adjusted;
this polynomial form remains available as a separate specialization. The bound also assumes $J\ge0$,
which holds for the source's $J$, a bound on norms: if there is any term,
$0\le\|h_k\|\le J$ implies it. The explicit assumption only adds a convention
for the empty family of terms, where the termwise norm bound is vacuous.
The shell-dependent comparison remains valid as stated. Its growth restriction
is no longer needed for the source's actual physical metric truncation:
\path{GaussianFilteredMetricTruncation.lean} proves this case by retaining the
distance of each omitted term in the Lieb--Robinson estimate. The graph
joining sites at distance at most $R$ has degree at most $S$, even when physical
metric balls have no useful cardinality bound. Choosing
$\mu=\log(S+1)+1$ makes the weighted omitted tail summable. With
$v=4S^2Je^{4\mu}$, $c_1=1/(4R(v+1))$ and
$\gamma=1/(2R)-vc_1>0$, the proved estimate has rate
$\tfrac12\min(c_1\Delta E,\gamma)>0$. Its prefactor depends only on
$R,S,J,\Delta E$, not on the finite site set or Hamiltonian. It applies for
$\ell\ge R$, the large-radius range of the source's order estimate, and
proves support, norm and Hermiticity for the actual $M_i$. The ground-space
and unique-ground-vector bounds and the estimate on $\|\sum_iM_i-H\|$ are
also proved. The zero-range case is exact and requires no division by $R$.
Explicit constants and the metric-to-graph comparison are recorded in
\href{https://sirui-lu.com/TNLean/paper-gaps/hastings05_geometric_growth_scope.pdf}
{the note on geometric growth}.

\section{The factorization length}
Immediately after eq.~(facld), the source says that $l_{fac}$ is the
minimum of $(c_1\Delta E)^{-1}$ and $\xi_C$. This reverses the elementary
rule for adding two decays: the sum is controlled by the slower decay,
hence by the larger length, or equivalently the smaller rate. The source's
following sentence, that $l_{fac}=(c_1\Delta E)^{-1}$ for small gap, agrees
with the larger-length reading. The explicit rate above and in
\eqref{eq:hastings05_truncated_dynamics_facld_decay} uses the minimum of
rates, with an additional factor two to absorb the time factor. This is a
local correction of the printed length comparison; the exponential order
asserted by eq.~(facld) is proved.
\bibliographystyle{unsrt}
\bibliography{references}
\end{document}
